\subsection{Overview of experimental requirements}

Experiments vary greatly in their requirements for kinematic range, supported event topology, and precision. Typically, the precision goals are informed by the size of other uncertainties, mainly the statistical precision, but also other unrelated systematic uncertainties, for example from geometry. But this can be misleading. While statistical uncertainty and often also other systematic sources have no or small correlation with kinematics, i.e.\ can be assumed to statistically average over many measurements, this is not true for radiative corrections, which have the potential to be highly correlated not only between data points of the same experiment, but even between multiple experiments if they use the same or similar formalism. 

Further, experiments might measure different event structures. In the case of $\ell p$ scattering, for example, the experiment could detect outgoing lepton, proton, or both, affecting the effectively seen corrections. Additionally, experiments can vary widely in the energy resolution, possibly requiring to describe the radiative tails accurately over rather large photon momenta ranges.  

\begin{table*}[t]
    \centering
    \caption{Characteristic parameters of published and upcoming proton form factor measurements}
      \scalebox{0.9}{
    \begin{tabular}{c|c|c|c|c|c|c}
         &\RotText{Measured particle}&\RotText{$E_\text{beam}$ [GeV]}&\RotText{$\theta$ [degrees]}&\RotText{$Q^2$ [GeV/$c^2$]}&\RotText{Effect of radius on form factor}&\RotText{Fractional contribution of $G_E$ to cross section}\\
         \hline
         GMp-12 \cite{Christy:2021snt}&$e^\prime$&2.222 -- 10.587&24.25 -- 53.5&1.577 -- 15.76&$>$100\%&0.2 -- 11.9\%\\

         MAMI FF \cite{A1:2010nsl,A1:2013fsc}&$e^\prime$&0.180 -- 0.855&15.5 -- 135&0.0033 -- 0.98&2 -- 59\%&0.23 -- 99.4\%\\
         MAMI High-$Q^2$ \cite{jmuellerthesis}&$e^\prime$&0.720 -- 1.500&15 -- 120&0.35 -- 1.95&$>$100\%&1.3 -- 92.5\%\\
         MAMI ISR \cite{Mihovilovic:2016rkr,Mihovilovic:2019jiz}&$e^\prime$&0.195, 0.330&15.21&0.001 -- 0.004&0.6 -- 2.4\%&$>$98\%\\
         MAMI Jet Target \cite{A1:2022wzx}&$e^\prime$&0.315&15 -- 40&0.007 -- 0.043&6 -- 39\%&71.2 -- 98.5\%\\
         PRad \cite{Xiong:2019umf}&$e^\prime$&1.1, 2.2&0.7 -- 6.5&0.00022 -- 0.058&0.13 -- 35\%&87.6 -- 100\%\\
        \hline
        AMBER \cite{Adams:2018pwt}&$\mu^\prime+p$&60, 100&N/A&0.001 -- 0.04&0.6 -- 24\%&91 -- 100\%\\
        MESA \cite{magixwebpage}&$e^\prime$&0.02 -- 0.105&15 -- 165&0.000027 -- 0.035
&0.016-21\%&91 -- 100\%\\
        MUSE \cite{MUSE:2017dod,Cline:2021ehf}&$e^\prime, \mu^\prime, \pi^\prime$&0.115 -- 0.21&20 -- 100&0.0016 -- 0.082
&0.96 -- 4.9\%&58 -- 99.6\%\\
        PRAD-II \cite{PRad:2020oor}&$e^\prime$&0.7, 1.4, 2.1&0.7 -- 6.5&0.00007 -- 0.056
&0.042 -- 34\%&88.6 -- 100\%\\
        ULQ2 \cite{Suda:2022hsm}&$e^\prime$&0.01 -- 0.06&0.7-6.5&0.000027 -- 0.012
&0.016 -- 7.2\%&56 -- 100\%\\


        
    \end{tabular}}

    
    \label{tab:expprop}
\end{table*}

Table \ref{tab:expprop} gives an overview over contemporary experiments \cite{Christy:2021snt,A1:2010nsl,A1:2013fsc,jmuellerthesis,Mihovilovic:2016rkr,Mihovilovic:2019jiz,A1:2022wzx,Xiong:2019umf,Adams:2018pwt,magixwebpage,MUSE:2017dod,Cline:2021ehf,PRad:2020oor,Suda:2022hsm}. We list beam energies, angle and $Q^2$ ranges, as well as detected particles. We also calculate an estimate of the effect of the proton radii on $G_E$ and $G_M$ as well as the fractional contribution of $G_E$ to the cross section, as an estimate for requirements on systematic uncertainties.  Of the planned experiments, only AMBER will measure the scattered lepton in coincidence with the recoiling proton. AMBER, PRad and PRAD-II both have forward kinematics with $\varepsilon$ very close to 1, while the other experiments tend to measure at larger angle ranges, $\varepsilon<1$. AMBER and MUSE will measure with leptons of both charges. The lower energy experiments require proper treatment of the lepton mass, and in the case of ULQ2, MUSE and likely MESA, a treatment without ultra-relativistic approximations.



Note that in most experiments at small $Q^2$, the absolute normalization of radiative corrections is less of an issue---it is generally left floating and determined by an extrapolating fit to $Q^2=0$. Typically experimental normalization uncertainties are significantly larger than radiative correction uncertainties, so that such a fit is unavoidable. On the other hand, experiments aimed to extract the radii typically measure at very small $Q^2$ to reduce systematical errors from the fit extrapolation to $Q^2=0$. Consequently, radiative corrections need to control relative, $Q^2$-dependent, errors to a very high degree: The proton radius puzzle is a 4\% difference in extracted radii, for recent reviews see Refs.~\cite{Antognini:2022xoo,Gao:2021sml,Peset:2021iul,Karr:2020wgh}. A meaningful measurement would extract the form factor slope to better than 0.5-1\% or so. A radius measurement at small $Q^2$ where the radius/slope effect is, say, 5\%, would then need to control $Q^2$-dependent systematic uncertainties over this $Q^2$ range to better than 0.05\%.  

For Rosenbluth separations at somewhat extreme kinematics, one of the form factor contributions to the measured cross section can be highly suppressed compared to the other one: $G_M$ is suppressed at small $Q^2$, while $G_E$ is at large $Q^2$. The uncertainty of radiative corrections on the cross section---in this case mainly the $\ep$-dependence at constant $Q^2$---then gets levered up by potentially large factors in the determination of the suppressed form factor. 

A further consideration is the requirement for unweighted events, or variance-reduced weights. This strongly depends on the overall speed of the remaining MC chain. For example, the Mainz A1 MC simulation code does not employ particle path tracing through magnetic fields and can simulate thousands if not millions of events per second, enabling to use fully weighted events, possibly even those with negative weights. In contrast, OLYMPUS' simulation chain included path tracing, full digitization and tracking for every simulated event, resulting in speeds of few events per second. This requires either unweighted events or at least measures to reduce the variance of the event weights. 






